\documentclass[11pt]{article}

\usepackage[margin=1in]{geometry}
\usepackage[T1]{fontenc}
\usepackage{amsmath,amssymb,mathtools,bm}
\usepackage{booktabs,longtable,array}
\usepackage{hyperref}

\hypersetup{
  colorlinks=true,
  linkcolor=blue,
  urlcolor=blue,
  citecolor=blue
}

\setlength{\parskip}{0.5em}
\setlength{\parindent}{0pt}
\allowdisplaybreaks

\newcommand{\avg}[1]{\left\langle #1 \right\rangle}
\newcommand{\dd}{\mathrm{d}}
\newcommand{\ii}{\mathrm{i}}
\newcommand{\kb}{k_{\mathrm{bin}}}
\newcommand{\kbm}{k_{\mathrm{bin},m}}
\newcommand{\cc}{c^{\mathrm{CC}}}
\newcommand{\rr}{\omega^{\mathrm{rfft}}}
\newcommand{\psihat}{\widehat{\psi}}
\newcommand{\qhat}{\widehat{q}}
\newcommand{\what}{\widehat{w}}
\newcommand{\thhat}{\widehat{\theta}}

\title{Reference for the Spectral Diagnostics Archive and Plots}
\author{}
\date{\today}

\begin{document}

\maketitle

\section{Purpose}

This note defines, in the notation actually used by the code, the arrays and
plots written under the archive
\begin{center}
\path{output_miquel_zero_tilt_galerkin_ars222_evolvemean_kebudget_blas1_Nx64_Nz256_dt5e5_t8/spectra/}.
\end{center}
The same definitions apply to the other ``\path{spectrum_history.npz}''
archives produced by the Chebyshev/Galerkin budget branch.

The point of this note is narrow: it explains what is stored and what is
plotted. It does \emph{not} attempt to re-derive the full NHQGE model from
first principles.

\section{State, transforms, and normalizations}

\subsection{Stored solver state}

The solver state is
\[
(\qhat,\what,\thhat,\Theta_{\mathrm{bar}}),
\]
with the following storage conventions:
\begin{itemize}
\item \(\qhat_n(k_x,k_y)\): Chebyshev coefficients of \(q^{\prime}\).
\item \(\what,\thhat\): Dirichlet Galerkin coefficients of \(w\) and
  \(\theta\), converted back to full Chebyshev coefficients before the
  diagnostics are evaluated.
\item \(\Theta_{\mathrm{bar}}\): Chebyshev coefficients of the mean-temperature
  deviation \(\Theta_{\mathrm{bar}}^{\prime}(z)\).
\end{itemize}

The streamfunction is recovered modewise by
\[
\psihat_n(k_x,k_y)
=
-\frac{\qhat_n(k_x,k_y)}{|k|^2 + L_d^{-2}}.
\]

\subsection{Vertical nodes and depth integration}

Let \(z_j\), \(j=0,\dots,N_z\), denote the Chebyshev--Gauss--Lobatto nodes on
\([0,1]\), and let \(\cc_j\) denote the Clenshaw--Curtis quadrature weights.
Whenever a diagnostic is depth-integrated in physical space, the code uses
\[
\int_0^1 f(z)\,\dd z
\approx
\sum_{j=0}^{N_z} \cc_j f(z_j).
\]

\subsection{Horizontal half-plane Parseval weights}

The horizontal fields are stored in the real-to-complex half plane. The code
therefore uses the standard \texttt{rfft2} Parseval weight
\[
\rr(k_x,k_y)
=
\begin{cases}
1, & k_y = 0 \text{ or Nyquist},\\
2, & \text{otherwise},
\end{cases}
\]
so that shell sums over the stored half-plane reproduce full-plane quadratic
quantities.

\subsection{Horizontal shell bins}

Let
\[
|k| = \sqrt{k_x^2 + k_y^2},
\qquad
\Delta k = \frac{2\pi}{L}.
\]
The archive stores shell centers
\[
\kbm = \left(m + \frac12\right)\Delta k,
\qquad
m=0,\dots,M-1,
\]
and shell \(m\) means
\[
m\Delta k \le |k| < (m+1)\Delta k.
\]
The array \path{k_bins} stores the centers \(\kbm\).

\subsection{A normalization warning}

The archive contains two different kinds of spectral quantities:
\begin{enumerate}
\item \textbf{physically normalized horizontal quadratic diagnostics},
  which include the factor \(N_x^{-4}\) and the \texttt{rfft2} weights;
\item \textbf{raw coefficient-amplitude diagnostics},
  which do \emph{not} include \(N_x^{-4}\), depth quadrature, or the
  \texttt{rfft2} weights.
\end{enumerate}

This distinction matters. In particular:
\begin{itemize}
\item \path{ke_horiz_spec} and all shell tendencies and cumulative fluxes are
  physically normalized.
\item \path{q_horiz_spec}, \path{w_horiz_spec}, \path{th_horiz_spec} are
  shell-binned but \emph{not} divided by \(N_x^4\).
\item \path{q_vert_spec}, \path{w_vert_spec}, \path{th_vert_spec} are raw
  Chebyshev coefficient norms.
\end{itemize}

\section{Scalar arrays in the archive}

The one-dimensional arrays \path{t} and \path{step} store the saved time and
time-step index.

The remaining one-dimensional scalar diagnostics are:
\begin{align*}
\mathrm{Nu}(t)
&=
1 + \avg{w\theta}_{xyz},\\
\avg{w\theta}_{xyz}
&=
\frac{1}{N_x^4}
\sum_{k_x,k_y}\rr(k_x,k_y)
\sum_{j=0}^{N_z}\cc_j
\Re\!\bigl(w_j(k_x,k_y)\,\theta_j(k_x,k_y)^*\bigr),\\
\mathrm{max\_speed}(t)
&=
\max_{j,x,y}\sqrt{u(z_j,x,y)^2 + v(z_j,x,y)^2},\\
\mathrm{max\_w}(t)
&=
\max_{j,x,y}|w(z_j,x,y)|,\\
\mathrm{max\_\theta}(t)
&=
\max_{j,x,y}|\theta(z_j,x,y)|,\\
\mathrm{max\_tw}(t)
&=
\max_{j,x,y}|w(z_j,x,y)\theta(z_j,x,y)|,
\end{align*}
with
\[
u = -\partial_y \psi,
\qquad
v = \partial_x \psi.
\]

The RMS quantities are
\begin{align*}
q_{\mathrm{rms}}^2
&=
\frac{1}{N_x^4}
\sum_{k_x,k_y}\rr
\sum_j \cc_j |q_j(k_x,k_y)|^2,\\
w_{\mathrm{rms}}^2
&=
\frac{1}{N_x^4}
\sum_{k_x,k_y}\rr
\sum_j \cc_j |w_j(k_x,k_y)|^2,\\
\theta_{\mathrm{rms}}^2
&=
\frac{1}{N_x^4}
\sum_{k_x,k_y}\rr
\sum_j \cc_j |\theta_j(k_x,k_y)|^2.
\end{align*}

The array \path{vol_avg_tw} stores \(\avg{w\theta}_{xyz}\) itself, and the
archive field \path{Nusselt} is defined by
\[
\mathrm{Nu} = 1 + \avg{w\theta}_{xyz}.
\]

The archive also stores the solver-consistent dealiased heat-flux profile used
by the prognostic mean-temperature equation. Let
\[
\avg{w\theta}_{xy}^{\mathrm{deal}}(z_j)
\]
denote the horizontal mean obtained from the same padded horizontal work grid
and inverse transforms used in the mean equation. Then
\begin{align*}
\avg{w\theta}_{xyz}^{\mathrm{deal}}
&=
\sum_{j=0}^{N_z}\cc_j \avg{w\theta}_{xy}^{\mathrm{deal}}(z_j),\\
\mathrm{Nu}^{\mathrm{deal}}
&=
1 + \avg{w\theta}_{xyz}^{\mathrm{deal}},\\
\Delta_{\mathrm{flux}}
&=
\avg{w\theta}_{xyz}^{\mathrm{deal}} - \avg{w\theta}_{xyz}.
\end{align*}
These are stored respectively under the keys \path{vol_avg_tw_dealiased},
\path{Nusselt_dealiased}, and \path{heat_flux_mismatch}.

The arrays \path{Nusselt_raw} and \path{vol_avg_tw_raw}, when present, are
simple aliases for \path{Nusselt} and \path{vol_avg_tw}; they are included only
to make the distinction between raw and dealiased thermal diagnostics explicit
in newer archives.

The array \path{th_bar_max} requires care. It is defined by
\[
\max_{0\le n\le N_z} |\Theta_{\mathrm{bar},n}(t)|.
\]
This is \emph{not} the physical-space maximum of
\(|\Theta_{\mathrm{bar}}^{\prime}(z)|\). It is the maximum absolute Chebyshev
coefficient of the mean-temperature deviation.

The mean-profile amplitudes and gradients are:
\begin{align*}
\Theta_{\max}
&=
\max_j |\Theta_{\mathrm{bar}}^{\prime}(z_j)|,\\
(\partial_z \Theta_{\mathrm{bar}}^{\prime})_{\max}
&=
\max_j |\partial_z \Theta_{\mathrm{bar}}^{\prime}(z_j)|,\\
g_{\min}
&=
\min_j \bigl(1 - \partial_z \Theta_{\mathrm{bar}}^{\prime}(z_j)\bigr),\\
g_{\max}
&=
\max_j \bigl(1 - \partial_z \Theta_{\mathrm{bar}}^{\prime}(z_j)\bigr),\\
g_{1/2}
&=
1 - \partial_z \Theta_{\mathrm{bar}}^{\prime}(z=1/2).
\end{align*}
These are stored as \path{th_bar_phys_max}, \path{dth_bar_dz_max},
\path{mean_grad_min}, \path{mean_grad_max}, and \path{mean_grad_mid}.

For \path{thermal_closure = "evolve_mean"}, the code also stores the global
mean-reservoir budget
\begin{align*}
\mathcal{E}_{\mathrm{mean}}
&=
\frac{1}{2\varepsilon^2}\int_0^1 \Theta_{\mathrm{bar}}^{\prime}(z)^2\,\dd z,\\
\mathcal{T}_{\mathrm{mean,flux}}
&=
\frac{1}{\varepsilon^2}\int_0^1
\Theta_{\mathrm{bar}}^{\prime}(z)\,
\partial_t\Theta_{\mathrm{bar}}^{\prime}\big|_{\mathrm{flux}}\,\dd z,\\
\mathcal{T}_{\mathrm{mean,diff}}
&=
\frac{1}{\varepsilon^2}\int_0^1
\Theta_{\mathrm{bar}}^{\prime}(z)\,
\partial_t\Theta_{\mathrm{bar}}^{\prime}\big|_{\mathrm{diff}}\,\dd z,\\
\mathcal{T}_{\mathrm{mean,tot}}
&=
\mathcal{T}_{\mathrm{mean,flux}}
+
\mathcal{T}_{\mathrm{mean,diff}}.
\end{align*}
These are stored under the keys \path{mean_energy},
\path{mean_flux_exchange_tendency}, \path{mean_diffusion_tendency}, and
\path{mean_total_tendency}.

The fluctuation-to-mean exchange residuals are stored in both raw and
solver-consistent dealiased form:
\begin{align*}
\mathcal{R}_{\mathrm{ex,raw}}
&=
\mathcal{T}_{\theta\leftrightarrow\mathrm{mean}}^{\mathrm{raw}}
+
\mathcal{T}_{\mathrm{mean,flux}},\\
\mathcal{R}_{\mathrm{ex,deal}}
&=
\mathcal{T}_{\theta\leftrightarrow\mathrm{mean}}^{\mathrm{deal}}
+
\mathcal{T}_{\mathrm{mean,flux}}.
\end{align*}
These correspond to the archive keys \path{mean_theta_exchange_residual} and
\path{mean_theta_exchange_residual_dealiased}; the fluctuation-side sums are
\path{th_mean_feedback_sum} and \path{th_mean_feedback_sum_dealiased}.
The corresponding \path{*_rel} fields divide these residuals by the sum of the
absolute values of the two compared transfer terms.

\section{Vertical and horizontal spectra}

\subsection{Vertical raw coefficient spectra}

Let \(\widetilde{w}_n(k_x,k_y)\) and \(\widetilde{\theta}_n(k_x,k_y)\) denote
the full Chebyshev coefficients obtained after lifting the stored Dirichlet
Galerkin coefficients back to the Chebyshev basis.

Then the arrays
\path{q_vert_spec}, \path{w_vert_spec}, and \path{th_vert_spec}
store
\begin{align*}
S^{(z)}_q(n)
&=
\sum_{k_x,k_y \in \mathrm{stored\ half\ plane}}
|\qhat_n(k_x,k_y)|^2,\\
S^{(z)}_w(n)
&=
\sum_{k_x,k_y \in \mathrm{stored\ half\ plane}}
|\widetilde{w}_n(k_x,k_y)|^2,\\
S^{(z)}_{\theta}(n)
&=
\sum_{k_x,k_y \in \mathrm{stored\ half\ plane}}
|\widetilde{\theta}_n(k_x,k_y)|^2.
\end{align*}

These are \emph{raw coefficient norms}. They do not include any depth
quadrature, any shell weighting, or any \(N_x^{-4}\) normalization.

The array \path{vert_mode} is simply the integer sequence
\[
(0,1,\dots,N_z).
\]

\subsection{Horizontal shell spectra of \(q\), \(w\), and \(\theta\)}

For a nodal spectral field \(f_j(k_x,k_y)\), define the shell-binned
depth-integrated amplitude
\[
S_f^{(k)}(m)
=
\sum_{|k|\in \mathrm{shell}\ m}
\rr(k_x,k_y)
\sum_{j=0}^{N_z}\cc_j |f_j(k_x,k_y)|^2.
\]

Then the archive stores \(S_q^{(k)}\), \(S_w^{(k)}\), and \(S_{\theta}^{(k)}\)
under the keys \path{q_horiz_spec}, \path{w_horiz_spec}, and
\path{th_horiz_spec}.

Again, these are shell amplitudes but not fully normalized physical energies,
because the code does \emph{not} divide them by \(N_x^4\).

\subsection{Horizontal kinetic-energy spectrum}

The stored shellwise kinetic-energy spectrum is
\[
E_{\mathrm{KE}}(m)
=
\sum_{|k|\in \mathrm{shell}\ m}
\frac{\rr(k_x,k_y)}{N_x^4}
\left[
\frac12 |k|^2
\sum_{j=0}^{N_z}\cc_j |\psi_j(k_x,k_y)|^2
\right].
\]
This quantity is stored as \path{ke_horiz_spec}.

\section{Split equations used in the shell budgets}

The code constructs the shell budgets from the same explicit/implicit split
used by the IMEX solver. At the diagnostic level, the relevant decomposition is
\begin{align*}
\partial_t q^{\prime}
&=
\mathcal{N}_q + \mathcal{B}_q + \mathcal{S}_q + \mathcal{D}_q,\\
\partial_t w
&=
\mathcal{N}_w + \mathcal{C}_{qw} + \mathcal{B}_w + \mathcal{D}_w,\\
\partial_t \theta
&=
\mathcal{N}_{\theta} + \mathcal{M}_{\theta} + \mathcal{C}_{\theta} + \mathcal{D}_{\theta},
\end{align*}
with
\begin{align*}
\mathcal{N}_q &= -A[\psi,q^{\prime}],&
\mathcal{B}_q &= -\ii \beta k_x \psihat,&
\mathcal{S}_q &= \partial_z w,&
\mathcal{D}_q &= -\lambda_q(|k|)\, q^{\prime},\\
\mathcal{N}_w &= -A[\psi,w],&
\mathcal{C}_{qw} &= \frac{1}{|k|^2 + L_d^{-2}}\,\partial_z q^{\prime},&
\mathcal{B}_w &= \frac{\widetilde{Ra}}{\sigma}\,\theta,&
\mathcal{D}_w &= -\lambda_w(|k|)\, w,\\
\mathcal{N}_{\theta} &= -A[\psi,\theta],&
\mathcal{M}_{\theta} &= -w\,\partial_z \Theta_{\mathrm{bar}}^{\prime},&
\mathcal{C}_{\theta} &= w,&
\mathcal{D}_{\theta} &= -\lambda_{\theta}(|k|)\,\theta.
\end{align*}

Here \(A[\psi,\cdot]\) denotes the dealiased horizontal Jacobian or flux-form
advection operator chosen by the run configuration, and
\(\lambda_q,\lambda_w,\lambda_{\theta}\) are the precomputed horizontal
dissipation multipliers stored internally as
\path{grid.diss_rate_q}, \path{grid.diss_rate_w}, and
\path{grid.diss_rate_th}.

For runs with \path{thermal_closure = "evolve_mean"}, the mean-temperature
equation used by the code is
\[
\partial_t \Theta_{\mathrm{bar}}^{\prime}
=
-\varepsilon^2 \partial_z \avg{w\theta}_{xy}
+ \frac{\varepsilon^2}{\sigma}\,\partial_{zz}\Theta_{\mathrm{bar}}^{\prime},
\]
but the shell budgets below only track how \(\Theta_{\mathrm{bar}}^{\prime}\)
feeds back into the \(\theta\) equation through \(\mathcal{M}_{\theta}\).

\section{Shellwise budget densities actually stored}

\subsection{Generic quadratic shell tendency}

For any nodal spectral field \(f_j(k_x,k_y)\) and nodal spectral term
\(g_j(k_x,k_y)\), the code defines the shell tendency
\[
\mathcal{T}[f;g](m)
=
\sum_{|k|\in \mathrm{shell}\ m}
\frac{\rr(k_x,k_y)}{N_x^4}
\sum_{j=0}^{N_z}\cc_j
\Re\!\bigl(f_j(k_x,k_y)^*\, g_j(k_x,k_y)\bigr).
\]
This is the shell contribution to
\(\frac{\dd}{\dd t}\bigl(\frac12 \|f\|_2^2\bigr)\).

\subsection{KE shell tendencies}

Because \(\psihat = -\qhat/(|k|^2+L_d^{-2})\), the kinetic-energy shell tendency
from a \(q\)-equation term \(q_t^{(\star)}\) is computed as
\[
\mathcal{T}^{\mathrm{KE}}_{\star}(m)
=
\sum_{|k|\in \mathrm{shell}\ m}
\frac{\rr(k_x,k_y)}{N_x^4}
\left[
-\frac{|k|^2}{|k|^2 + L_d^{-2}}
\sum_{j=0}^{N_z}\cc_j
\Re\!\bigl(\psi_j(k_x,k_y)^*\, q_{t,j}^{(\star)}(k_x,k_y)\bigr)
\right].
\]

The archive stores the following KE shell tendencies:
\begin{align*}
T^{\mathrm{KE}}_{\mathrm{NL}} &= \mathcal{T}^{\mathrm{KE}}_{\mathcal{N}_q},\\
T^{\mathrm{KE}}_{\beta} &= \mathcal{T}^{\mathrm{KE}}_{\mathcal{B}_q},\\
T^{\mathrm{KE}}_{\mathrm{stretch}} &= \mathcal{T}^{\mathrm{KE}}_{\mathcal{S}_q},\\
T^{\mathrm{KE}}_{\mathrm{diss}} &= \mathcal{T}^{\mathrm{KE}}_{\mathcal{D}_q},\\
T^{\mathrm{KE}}_{\mathrm{tot}} &=
T^{\mathrm{KE}}_{\mathrm{NL}}
+ T^{\mathrm{KE}}_{\beta}
+ T^{\mathrm{KE}}_{\mathrm{stretch}}
+ T^{\mathrm{KE}}_{\mathrm{diss}}.
\end{align*}
These correspond respectively to the archive keys
\path{ke_nonlinear_shell_tendency},
\path{ke_beta_shell_tendency},
\path{ke_stretch_shell_tendency},
\path{ke_diss_shell_tendency}, and
\path{ke_total_shell_tendency}.

The nonlinear cumulative flux is defined shellwise by
\[
\Pi^{\mathrm{KE}}_{\mathrm{NL}}(m)
=
-\sum_{\ell=0}^{m}
\mathcal{T}^{\mathrm{KE}}_{\mathcal{N}_q}(\ell),
\]
and is stored as \path{ke_nonlinear_flux}.

\subsection{\(w\)-variance shell tendencies}

The \(w\)-variance shell budgets use the generic quadratic form
\(\mathcal{T}[w; \cdot]\). The stored quantities are
\begin{align*}
T^{w}_{\mathrm{NL}} &= \mathcal{T}[w;\mathcal{N}_w],\\
T^{w}_{q} &= \mathcal{T}[w;\mathcal{C}_{qw}],\\
T^{w}_{\mathrm{buoy}} &= \mathcal{T}[w;\mathcal{B}_w],\\
T^{w}_{\mathrm{diss}} &= \mathcal{T}[w;\mathcal{D}_w],\\
T^{w}_{\mathrm{tot}} &= T^{w}_{\mathrm{NL}} + T^{w}_{q}
+ T^{w}_{\mathrm{buoy}} + T^{w}_{\mathrm{diss}}.
\end{align*}
These correspond to the archive keys
\path{w_nonlinear_shell_tendency},
\path{w_q_coupling_shell_tendency},
\path{w_buoyancy_shell_tendency},
\path{w_diss_shell_tendency}, and
\path{w_total_shell_tendency}.

The nonlinear cumulative flux is
\[
\Pi^{w}_{\mathrm{NL}}(m)
=
-\sum_{\ell=0}^{m}
\mathcal{T}[w;\mathcal{N}_w](\ell),
\]
stored as \path{w_nonlinear_flux}.

\subsection{\(\theta\)-variance shell tendencies}

The \(\theta\)-variance shell budgets use the same quadratic form:
\begin{align*}
T^{\theta}_{\mathrm{NL}} &= \mathcal{T}[\theta;\mathcal{N}_{\theta}],\\
T^{\theta}_{\mathrm{mf}} &= \mathcal{T}[\theta;\mathcal{M}_{\theta}],\\
T^{\theta}_{\mathrm{cond}} &= \mathcal{T}[\theta;\mathcal{C}_{\theta}],\\
T^{\theta}_{\mathrm{diss}} &= \mathcal{T}[\theta;\mathcal{D}_{\theta}],\\
T^{\theta}_{\mathrm{tot}} &= T^{\theta}_{\mathrm{NL}} + T^{\theta}_{\mathrm{mf}}
+ T^{\theta}_{\mathrm{cond}} + T^{\theta}_{\mathrm{diss}}.
\end{align*}
These correspond to the archive keys
\path{th_nonlinear_shell_tendency},
\path{th_mean_feedback_shell_tendency},
\path{th_conduction_shell_tendency},
\path{th_diss_shell_tendency}, and
\path{th_total_shell_tendency}.

The array \path{th_conduction_shell_tendency} deserves a short comment. In the
split code it comes from the implicit source \(+\;w\) in the \(\theta\)
equation, so the shell diagnostic is
\[
\mathcal{T}[\theta; w](m).
\]
The label ``conduction'' is therefore an interpretation of the term's origin in
the conductive background-gradient forcing, not a separate Laplacian acting on
\(\theta\).

The nonlinear cumulative flux is
\[
\Pi^{\theta}_{\mathrm{NL}}(m)
=
-\sum_{\ell=0}^{m}
\mathcal{T}[\theta;\mathcal{N}_{\theta}](\ell),
\]
stored as \path{th_nonlinear_flux}.

\subsection{Dealiased thermal shell diagnostics}

The archive also stores shellwise thermal quantities built from the same
dealiased horizontal product path used by the mean equation. Let
\[
F_m^{\mathrm{deal}}(z_j)
=
\avg{w_m\theta_m}_{xy}^{\mathrm{deal}}(z_j)
\]
be the dealiased shell-filtered mean heat-flux profile. Then the stored fields
are
\begin{align*}
H_m^{\mathrm{deal}}
&=
\sum_{j=0}^{N_z}\cc_j F_m^{\mathrm{deal}}(z_j),\\
\mathcal{T}^{\theta,\mathrm{deal}}_{\mathrm{cond}}(m)
&=
H_m^{\mathrm{deal}},\\
\mathcal{T}^{w,\mathrm{deal}}_{\mathrm{buoy}}(m)
&=
\frac{\widetilde{Ra}}{\sigma}\,
H_m^{\mathrm{deal}},\\
\mathcal{T}^{\theta,\mathrm{deal}}_{\mathrm{mf}}(m)
&=
-\sum_{j=0}^{N_z}\cc_j\,
\partial_z\Theta_{\mathrm{bar}}^{\prime}(z_j)\,
F_m^{\mathrm{deal}}(z_j).
\end{align*}
These are stored as \path{heat_flux_shell_dealiased},
\path{th_conduction_shell_tendency_dealiased},
\path{w_buoyancy_shell_tendency_dealiased}, and
\path{th_mean_feedback_shell_tendency_dealiased}.

These are not independent new PDE terms; they are the shellwise thermal
exchange diagnostics obtained after replacing the raw product \(w\theta\) by the
same dealiased product path used in the mean equation itself.

\section{What each PNG in \texttt{spectra/} means}

\subsection{The two spectrum heatmaps}

The run script writes two immediate summary figures:
\begin{itemize}
\item \path{w_vertical_spectrum_vs_time.png}
\item \path{w_horizontal_spectrum_vs_time.png}
\end{itemize}

These are heatmaps of
\[
\log_{10}\!\bigl(\max(S,10^{-300})\bigr),
\]
where \(S\) is respectively \path{w_vert_spec} or \path{w_horiz_spec}. The
color scale is linear in this logged quantity, using the minimum and maximum
values present in the saved archive.

So these two PNGs are purely visualization transforms of the stored arrays:
\begin{itemize}
\item \path{w_vertical_spectrum_vs_time.png} plots
  \(\log_{10}\!\bigl(\max(S_w^{(z)},10^{-300})\bigr)\).
\item \path{w_horizontal_spectrum_vs_time.png} plots
  \(\log_{10}\!\bigl(\max(S_w^{(k)},10^{-300})\bigr)\).
\end{itemize}

\subsection{Budget heatmaps and lineouts}

The directory \path{ke_budget_plots/} contains heatmaps and lineouts for the
KE, \(w\)-variance, and \(\theta\)-variance shell budgets.

If \(X(t,m)\) is any stored signed shell quantity such as
\path{ke_stretch_shell_tendency} or \path{th_mean_feedback_shell_tendency}, the
plotter displays
\[
\mathrm{symlog}(X)
=
\mathrm{sgn}(X)\,
\log_{10}\!\left(1 + \frac{|X|}{\Lambda}\right),
\]
where the threshold \(\Lambda\) is chosen automatically from the data percentiles.

This transform is \emph{only} for plotting. The archive stores the unmodified
arrays.

Therefore:
\begin{itemize}
\item each \path{*_heatmap.png} in \path{ke_budget_plots/} is a heatmap of the
  corresponding stored shell field \(X(t,m)\) after this signed symlog
  transform;
\item each \path{*_lineouts.png} is a set of fixed-time slices
  \(m \mapsto X(t_{\mathrm{sel}},m)\), again plotted after the same signed
  symlog transform.
\end{itemize}

\subsection{Mean-exchange history plot}

The directory \path{mean_exchange_plots/} contains
\path{mean_exchange_history.png}. This figure is a collection of time traces of
the stored scalar arrays listed above:
\begin{itemize}
\item mean-reservoir budget terms:
  \path{mean_flux_exchange_tendency},
  \path{th_mean_feedback_sum},
  \path{th_mean_feedback_sum_dealiased},
  \path{mean_diffusion_tendency},
  \path{mean_theta_exchange_residual},
  \path{mean_theta_exchange_residual_dealiased},
  \path{mean_total_tendency};
\item heat-flux diagnostics:
  \path{Nusselt_dealiased},
  \path{Nusselt},
  \path{vol_avg_tw_dealiased},
  \path{vol_avg_tw},
  \path{heat_flux_mismatch};
\item mean-gradient and mean-amplitude diagnostics:
  \path{mean_grad_min}, \path{mean_grad_mid}, \path{mean_grad_max},
  \path{th_bar_phys_max}, \path{dth_bar_dz_max}, \path{mean_energy}.
\end{itemize}
The first, second, and fourth panels use the same signed symlog display
transform as the shell-budget plots; the gradient panel is plotted on a linear
axis.

\section{Dominant-shell products}

The directory \path{dominant_shell/} is generated from the saved archive as
follows.

\subsection{Shell selection}

Choose a selector field \(X_{\mathrm{sel}}(t,m)\). By default the script uses
the stored KE stretching field:
\[
X_{\mathrm{sel}}(t,m) = T^{\mathrm{KE}}_{\mathrm{stretch}}(t,m).
\]

Choose a reference time \(t_{\ast}\). By default the script uses the final saved
time.

Then the dominant shell index is
\[
m_{\ast}
=
\arg\max_m |X_{\mathrm{sel}}(t_{\ast},m)|.
\]
The archive stores both \(m_{\ast}\) and the corresponding shell center
\(k_{\mathrm{bin},m_{\ast}}\) in
\path{dominant_shell_timeseries.npz}.

\subsection{Extracted time series}

At this fixed shell \(m_{\ast}\), the script extracts the time series
\[
t \mapsto X(t,m_{\ast})
\]
for the following quantities:
\begin{itemize}
\item KE: \path{ke_total}, \path{ke_stretch}, \path{ke_diss}, \path{ke_nonlinear}.
\item \(w\)-variance: \path{w_total}, \path{w_buoyancy}, \path{w_q_coupling},
  \path{w_diss}, \path{w_nonlinear}.
\item \(\theta\)-variance: \path{th_total}, \path{th_conduction},
  \path{th_mean_feedback}, \path{th_diss}, \path{th_nonlinear}.
\end{itemize}
together with the scalar time traces \path{Nusselt} and \path{max_speed}.

The figure \path{dominant_shell_timeseries.png} plots these extracted shellwise
time series using the same signed symlog display convention as the budget
lineouts.

\section{Key archive fields at a glance}

\begin{longtable}{>{\raggedright\arraybackslash}p{0.30\textwidth} >{\raggedright\arraybackslash}p{0.62\textwidth}}
\toprule
Field & Meaning \\
\midrule
\endfirsthead
\toprule
Field & Meaning \\
\midrule
\endhead
\path{t}, \path{step} & Saved times and time-step indices. \\
\path{Nusselt}, \path{Nusselt_raw} & Raw \(1 + \avg{w\theta}_{xyz}\). \\
\path{Nusselt_dealiased} & Solver-consistent \(1 + \avg{w\theta}_{xy}^{\mathrm{deal}}\). \\
\path{vol_avg_tw}, \path{vol_avg_tw_raw} & Raw volume-averaged convective flux \(\avg{w\theta}_{xyz}\). \\
\path{vol_avg_tw_dealiased} & Solver-consistent dealiased volume-averaged convective flux. \\
\path{heat_flux_mismatch} & Difference between the dealiased and raw volume-averaged heat flux. \\
\path{max_speed} & \(\max \sqrt{u^2+v^2}\) on the saved nodal/physical grid. \\
\path{max_w}, \path{max_theta}, \path{max_tw} & Maxima of \(|w|\), \(|\theta|\), and \(|w\theta|\) on the saved nodal/physical grid. \\
\path{q_rms}, \path{w_rms}, \path{th_rms} & RMS values from depth-integrated, horizontally Parseval-weighted quadratic norms. \\
\path{th_bar_max} & Maximum absolute Chebyshev coefficient of \(\Theta_{\mathrm{bar}}^{\prime}\), not a physical-space maximum. \\
\path{th_bar_phys_max}, \path{dth_bar_dz_max} & Physical-space maxima of \(|\Theta_{\mathrm{bar}}^{\prime}|\) and \(|\partial_z\Theta_{\mathrm{bar}}^{\prime}|\). \\
\path{mean_grad_min}, \path{mean_grad_mid}, \path{mean_grad_max} & Min, midplane value, and max of \(1-\partial_z\Theta_{\mathrm{bar}}^{\prime}\). \\
\path{mean_energy}, \path{mean_flux_exchange_tendency}, \path{mean_diffusion_tendency}, \path{mean_total_tendency} & Global mean-reservoir quadratic energy and its split tendencies. \\
\path{th_mean_feedback_sum}, \path{th_mean_feedback_sum_dealiased} & Raw and dealiased fluctuation-side thermal exchange sums. \\
\path{mean_theta_exchange_residual}, \path{mean_theta_exchange_residual_dealiased} & Raw and dealiased fluctuation/mean exchange residuals. \\
\path{vert_mode} & Integer Chebyshev mode index \(n=0,\dots,N_z\). \\
\path{q_vert_spec}, \path{w_vert_spec}, \path{th_vert_spec} & Raw vertical Chebyshev coefficient norms. \\
\path{k_bins} & Horizontal shell centers \(\kbm=(m+\tfrac12)\Delta k\). \\
\path{q_horiz_spec}, \path{w_horiz_spec}, \path{th_horiz_spec} & Depth-integrated shell amplitudes of \(q\), \(w\), \(\theta\), without \(N_x^{-4}\) normalization. \\
\path{ke_horiz_spec} & Physically normalized shellwise kinetic-energy spectrum. \\
\path{ke_*_shell_tendency} & Signed shell contributions to \(\dd(\tfrac12\|\nabla_h\psi\|^2)/\dd t\). \\
\path{w_*_shell_tendency} & Signed shell contributions to \(\dd(\tfrac12\|w\|^2)/\dd t\). \\
\path{th_*_shell_tendency} & Raw signed shell contributions to \(\dd(\tfrac12\|\theta\|^2)/\dd t\). \\
\path{heat_flux_shell_dealiased}, \path{th_*_shell_tendency_dealiased}, \path{w_buoyancy_shell_tendency_dealiased} & Solver-consistent dealiased thermal shell diagnostics built from the mean-equation heat-flux path. \\
\path{ke_nonlinear_flux}, \path{w_nonlinear_flux}, \path{th_nonlinear_flux} & Cumulative shellwise nonlinear fluxes, defined as minus the cumulative sum of the corresponding nonlinear shell tendency. \\
\bottomrule
\end{longtable}

\section{Practical reading guide}

For the current budget archives, the cleanest interpretation route is:
\begin{enumerate}
\item use \path{w_vertical_spectrum_vs_time.png} and
  \path{w_horizontal_spectrum_vs_time.png} for a quick view of where amplitude
  is moving;
\item use the \path{ke_budget_plots/*_heatmap.png} files to identify which
  shellwise source terms are growing, at which times and horizontal shells;
\item use \path{dominant_shell/dominant_shell_timeseries.png} and
  \path{dominant_shell/dominant_shell_summary.md} to inspect the exact
  low-shell pathway at the shell selected by the dominant KE stretching signal.
\end{enumerate}

The important distinction is that the spectrum PNGs are nonnegative log-heatmaps
of stored amplitudes, whereas the budget PNGs are signed symlog visualizations
of shellwise source/sink terms.

\end{document}
